\documentclass[11pt]{article}
% ---------- Packages ----------
\usepackage[margin=1in]{geometry}   % 1-inch page margins
\usepackage{graphicx}               % lets us insert figures (PDFs)
\usepackage{booktabs}               % clean table lines (\toprule, \midrule)
\usepackage{caption}                % nicer captions
\usepackage{subcaption}             % several plots side by side in one figure
\usepackage{float}                  % [H] keeps figures where we put them
\usepackage[hidelinks]{hyperref}    % clickable links (GitHub URL)

\captionsetup{font=small, labelfont=bf}  % "Figure 1." in bold, small caption text

% ---------- Title ----------
\title{Expression of \textit{IL24}, \textit{ABCA9} and \textit{ABCB7}
in Relation to Maternal Age, BMI and Preeclampsia}
\author{Jubilee Mwasa}
\date{November 11, 2026}

\begin{document}
\maketitle

% ================= METHODS =================
\section{Methods}
\subsection{Data Source}
Data were obtained from a nested case--control gene expression study within
the Vitamin D Antenatal Asthma Reduction Trial (VDAART) \cite{mirzakhani2016},
available from the Gene Expression Omnibus (accession GSE85307). Peripheral
blood gene expression was measured at 10--18 weeks of gestation using the
Affymetrix Human Gene 1.0 ST Array. The initial dataset contained 157 samples.

\subsection{Data Processing and Sample Size}
Expression samples were matched to participant information using the GEO
sample identifier (geoID), and all 157 samples were successfully matched.
Probe IDs were mapped to gene symbols, and expression values from multiple
probes for the same gene were averaged. Covariate values were converted from
text to numeric or categorical form. Maternal BMI was categorized using World
Health Organization cut-points (underweight $<18.5$, normal weight
18.5--24.9, overweight 25--29.9, obesity $\geq 30$~kg/m$^2$). No participants
had missing data for pregnancy outcome, maternal age or BMI, so no samples
were removed and the final sample included 157 participants (110 normal
pregnancies and 47 preeclampsia cases).

\subsection{Statistical Analysis}
Three genes were selected for analysis: \textit{IL24}, \textit{ABCA9} and
\textit{ABCB7}. For each gene, a histogram of expression, a scatterplot of
expression against maternal age with a LOESS trend line, and a boxplot of
expression by pregnancy outcome and BMI category were generated. Associations
were assessed with Spearman correlation (maternal age), Wilcoxon rank-sum
tests (pregnancy outcome) and Kruskal--Wallis tests (BMI category), with
$p < 0.05$ considered statistically significant. Participant characteristics
were summarized as mean (SD) for continuous variables and $n$ (\%) for
categorical variables, stratified by pregnancy outcome. A heatmap of ten genes
was generated using row-scaled (z-score) expression with hierarchical
clustering of genes and participants (Euclidean distance, complete linkage).
Finally, a dumbbell plot was used to compare mean z-scored expression of the
same ten genes between pregnancy outcome groups.

\subsection{Software}
All analyses were performed in R version 4.6.1 \cite{R-base}. Data were
processed using dplyr (v1.2.1) \cite{R-dplyr} and tidyr (v1.3.2)
\cite{R-tidyr}. Figures were created with ggplot2 (v4.0.3) \cite{R-ggplot2}
and pheatmap (v1.0.13) \cite{R-pheatmap}, and the report was generated with
knitr (v1.51) \cite{R-knitr}.

\subsection{Code Availability}
All code is available at \url{https://github.com/YOUR-USERNAME/YOUR-REPO}.

\section{Results}
\subsection{Participant Characteristics}
Maternal age was similar between groups, but women who developed
preeclampsia had lower vitamin D levels, delivered earlier, and were more
often classified as obese than women with normal pregnancies
(Table~\ref{tab:summary}).

\begin{table}[H]
\centering
\caption{Participant characteristics stratified by pregnancy outcome.
Continuous variables are presented as mean (SD) and categorical variables
as $n$ (\%). Vitamin D was measured as serum 25-hydroxyvitamin D at baseline.}
\label{tab:summary}
\small
\input{table1.tex}
\end{table}

\subsection{Distribution of Gene Expression}
  Expression of \textit{IL24}, \textit{ABCA9} and \textit{ABCB7} was unimodal
  and approximately symmetric, with a slight right tail for \textit{IL24} and
  \textit{ABCA9} (Figure~\ref{fig:hist}).

\begin{figure}[H]
\centering
\begin{subfigure}{0.32\textwidth}
  \centering
  \includegraphics[width=\linewidth]{figures/IL24_histogram.pdf}
  \caption{\textit{IL24}}
\end{subfigure}
\begin{subfigure}{0.32\textwidth}
  \centering
  \includegraphics[width=\linewidth]{figures/ABCA9_histogram.pdf}
  \caption{\textit{ABCA9}}
\end{subfigure}
\begin{subfigure}{0.32\textwidth}
  \centering
  \includegraphics[width=\linewidth]{figures/ABCB7_histogram.pdf}
  \caption{\textit{ABCB7}}
\end{subfigure}
\caption{Distribution of expression of (a) \textit{IL24}, (b) \textit{ABCA9}
and (c) \textit{ABCB7} across all 157 participants.}
\label{fig:hist}
\end{figure}

\subsection{Gene Expression and Maternal Age}
None of the three genes was significantly associated with maternal age.
\textit{IL24} ($\rho = 0.12$, $p = 0.14$) and \textit{ABCA9}
($\rho = -0.04$, $p = 0.60$) showed no meaningful relationship, while
\textit{ABCB7} showed a weak, non-significant negative trend
($\rho = -0.15$, $p = 0.057$) (Figure~\ref{fig:scatter}).

\begin{figure}[H]
\centering
\begin{subfigure}{0.32\textwidth}
  \centering
  \includegraphics[width=\linewidth]{figures/IL24_scatter.pdf}
  \caption{\textit{IL24}}
\end{subfigure}
\begin{subfigure}{0.32\textwidth}
  \centering
  \includegraphics[width=\linewidth]{figures/ABCA9_scatter.pdf}
  \caption{\textit{ABCA9}}
\end{subfigure}
\begin{subfigure}{0.32\textwidth}
  \centering
  \includegraphics[width=\linewidth]{figures/ABCB7_scatter.pdf}
  \caption{\textit{ABCB7}}
\end{subfigure}
\caption{Gene expression versus maternal age for (a) \textit{IL24},
(b) \textit{ABCA9} and (c) \textit{ABCB7}. Each point represents one
participant; the line shows a LOESS fit with a 95\% confidence band.}
\label{fig:scatter}
\end{figure}

\subsection{Gene Expression by Pregnancy Outcome and BMI}
Expression of \textit{IL24}, \textit{ABCA9} and \textit{ABCB7} did not
differ between normal and preeclamptic pregnancies ($p = 0.60$, $0.43$
and $0.81$, respectively) or across maternal BMI categories ($p = 0.78$,
$0.49$ and $0.15$, respectively) (Figure~\ref{fig:box}).

\begin{figure}[H]
\centering
\begin{subfigure}{0.49\textwidth}
  \centering
  \includegraphics[width=\linewidth]{figures/IL24_boxplot.pdf}
  \caption{\textit{IL24}}
\end{subfigure}
\begin{subfigure}{0.49\textwidth}
  \centering
  \includegraphics[width=\linewidth]{figures/ABCA9_boxplot.pdf}
  \caption{\textit{ABCA9}}
\end{subfigure}

\vspace{0.5em}
\begin{subfigure}{0.49\textwidth}
  \centering
  \includegraphics[width=\linewidth]{figures/ABCB7_boxplot.pdf}
  \caption{\textit{ABCB7}}
\end{subfigure}
\caption{Gene expression by pregnancy outcome and maternal BMI category for
(a) \textit{IL24}, (b) \textit{ABCA9} and (c) \textit{ABCB7}. Boxes show the
median and interquartile range; points represent individual participants.}
\label{fig:box}
\end{figure}

\subsection{Clustering of Ten Genes}
Hierarchical clustering of ten genes did not separate participants by
pregnancy outcome or maternal BMI category, as both tracking bars were
intermixed across clusters (Figure~\ref{fig:heatmap}).

\begin{figure}[H]
\centering
\includegraphics[width=0.95\textwidth]{figures/heatmap.pdf}
\caption{Heatmap of row-scaled (z-score) expression for ten genes across
157 participants. Genes (rows) and participants (columns) were clustered
hierarchically using Euclidean distance and complete linkage. Tracking bars
indicate pregnancy outcome and maternal BMI category. Red indicates higher
and blue lower expression relative to each gene's mean; colors are capped
at $\pm 3$ SD.}
\label{fig:heatmap}
\end{figure}


\subsection{Differences in Mean Expression Between Outcome Groups}
Differences in mean expression between outcome groups were small for all
ten genes (less than 0.4 standard deviations), and the 95\% confidence
intervals of the two groups overlapped for every gene. \textit{ATG101},
\textit{ANKRD60} and \textit{AKR1C8} showed the largest increases in
preeclampsia, whereas \textit{CATSPERT} and \textit{BHMT} were slightly
lower (Figure~\ref{fig:dumbbell}).

\begin{figure}[H]
\centering
\includegraphics[width=0.75\textwidth]{figures/dumbbell.pdf}
\caption{Dumbbell plot of mean z-scored expression for ten genes in normal
($n = 110$) and preeclamptic ($n = 47$) pregnancies. Points show group means
with 95\% confidence intervals; grey bars connect the two group means for
each gene. Genes are ordered by the difference between groups.}
\label{fig:dumbbell}
\end{figure}
\bibliographystyle{unsrt}     % numbers references in the order they are cited
\bibliography{references}     % uses references.bib

\end{document}
